\documentclass[11pt]{article}

\usepackage[margin=1in]{geometry}
\usepackage{amsfonts,amsmath}
\usepackage{booktabs}
\usepackage{multirow}
\usepackage{graphicx}
\usepackage{xcolor}
\usepackage[colorlinks,citecolor=blue,urlcolor=blue,linkcolor=blue]{hyperref}
\usepackage[numbers,sort&compress]{natbib}
\usepackage[affil-it]{authblk}
\usepackage{lineno}

\renewcommand{\arraystretch}{1.3}
\emergencystretch=1.5em

\newcommand{\Gr}{\mathrm{Gr}}
\newcommand{\R}{\mathbb{R}}

\title{Predictive Validity Without Discriminative Validity:\\
Cell-Type Centroid Metrics Score Trained Embeddings Below Random}

\author{Elliot Tower}
\affil{Independent Researcher \\ \texttt{elliot@elliottower.ai}}
\date{}

\begin{document}
\linenumbers
\maketitle

\begin{abstract}
Single-cell foundation models are ranked using embedding evaluation
metrics, but whether high scores reflect biological signal or geometric
artifacts of high-dimensional spaces is unknown. We evaluated twelve
metrics---five scIB bio-conservation, two centroid-alignment, and five
geometric---on 104 embedding conditions across 25 tissues, using
cross-technology cell-type transfer as the downstream criterion (10x
Chromium 3$'$ v3 $\rightarrow$ Smart-seq2). Procrustes similarity on
cell-type centroids achieves the strongest prediction (partial
$\rho = +0.66$, $p < 10^{-13}$), followed by silhouette and CKA
($\rho = +0.54$, $p < 10^{-8}$). CKA generalizes to two of three
independent datasets (pancreas $\rho = +0.59$, Tabula Sapiens
$\rho = +0.60$). All 104 conditions score below random Gaussian
centroid matrices on both CKA (real mean 0.74 vs.\ null 0.92--0.99)
and Procrustes (real mean 0.78 vs.\ null 0.96--0.99). The metrics
rank real embeddings correctly relative to each other while placing
every trained model below unstructured noise. Bio-conservation
metrics respond non-monotonically to controlled embedding degradation
(17--22 of 24 conditions). These results establish that predictive
validity---correct relative ranking---and discriminative
validity---distinguishing learned from random by absolute score---are
orthogonal properties in high-dimensional embedding evaluation.
Relative rankings remain reliable; absolute thresholds do not.
\end{abstract}

\noindent\textbf{Keywords:} single-cell foundation models, embedding evaluation, construct validity, centered kernel alignment, Procrustes analysis, scIB

\section{Introduction}
\label{sec:intro}

Single-cell foundation model benchmarks consistently find that PCA on
raw counts matches or exceeds large pretrained models across most
evaluation axes \citep{wu2025scfm, vcbench2026, unifiedscfm2026,
cellxgeneintegration2026}. \citet{wu2025scfm} raised the possibility
that the metrics themselves might be at fault. The question is whether
high metric scores reflect biological signal learned by the model, or
geometric properties that arise in any high-dimensional embedding.

Two distinct validity properties bear on this question. \emph{Predictive
validity} asks whether metric scores correlate with downstream
performance: does a higher-scoring embedding transfer better?
\emph{Discriminative validity} asks whether the metric can distinguish
trained from random embeddings by absolute score: does crossing a
threshold guarantee biological content? These properties are logically
independent. A metric can rank trained models correctly (predictive)
while scoring all of them below random noise (non-discriminative), or
it can separate trained from random (discriminative) while failing to
rank among trained models (non-predictive). Treating one as evidence for
the other conflates two distinct claims about metric quality.

We evaluate twelve metrics spanning three families---scIB
bio-conservation \citep{luecken2022scib}, cell-type centroid alignment,
and five geometric modules---on 104 embedding conditions across 25
tissues. Ground truth is cross-technology cell-type transfer F1 from
10x Chromium 3$'$ v3 to Smart-seq2.

Three findings emerge. First, most metrics predict transfer: Procrustes
similarity on cell-type centroids achieves partial $\rho = +0.66$, and
nine of twelve metrics reach significance after controlling for
dimensionality. Second, all 104 conditions score below random Gaussian
centroid matrices on both CKA and Procrustes---the metrics predict
relative quality but cannot certify that an embedding contains
biological signal. Third, bio-conservation metrics respond
non-monotonically to controlled embedding degradation, with scores
increasing at intermediate noise levels before collapsing.

These findings motivate a minimum standard for embedding evaluation:
report relative rankings with matched-dimensionality baselines rather
than absolute scores, and validate metric behavior under controlled
degradation.


\section{Results}

\subsection{Evaluation panel}
\label{sec:panel}

We constructed a cross-technology transfer panel from CellxGene Census
(v2023-12-15). Source embeddings are computed on 10x Chromium 3$'$ v3
cells; target labels come from Smart-seq2 (EFO:0008931) cells in the
same tissue with zero donor overlap. Tissues were auto-discovered by
requiring $\geq 200$ cells per assay and $\geq 8$ shared cell types,
yielding 25 tissues. Eight embedding models
span three architectural families: Geneformer \citep{theodoris2023geneformer}
($d = 512$), Geneformer-v2 104M and 316M, scGPT \citep{cui2024scgpt}
($d = 512$), scVI \citep{lopez2018scvi} ($d = 50$), bag-of-genes
PCA-512 ($d = 512$; following \citealp{cole2026morpheus}), random
Gaussian projection ($d = 512$), and an untrained two-layer MLP encoder
($d = 512$). Six contender models (excluding the two null baselines)
produce 104 contender conditions. Ground truth is macro-averaged
logistic-regression F1 on shared cell types, trained on source
embeddings and evaluated on target.


\subsection{Metric prediction of cross-technology transfer}
\label{sec:prediction}

Twelve metrics were evaluated on all 104 contender conditions. Partial
Spearman rank correlations control for embedding dimensionality via rank
residuals from OLS; tissue-stratified permutation tests (10,000
iterations) provide $p$-values robust to within-tissue dependence.

Nine of twelve metrics significantly predict cross-technology transfer
(Table~\ref{tab:metrics}). Procrustes similarity on cell-type centroids
achieves the strongest correlation (partial $\rho = +0.66$,
$p < 10^{-13}$), followed by target-domain silhouette ($\rho = +0.65$)
and silhouette on cell-type labels ($\rho = +0.60$). Cell-type CKA
achieves $\rho = +0.53$ ($p < 10^{-8}$). Among scIB metrics, graph
connectivity ($\rho = +0.58$), NMI ($\rho = +0.55$), and kNN purity
($\rho = +0.52$) predict transfer at comparable strength. ARI predicts
more weakly ($\rho = +0.37$). The three failures are
direction stability ($\rho = +0.02$, $p = 0.83$) and
participation ratio ($\rho = +0.45$, nominally significant but
confounded by dimensionality; Supplementary~\ref{app:modules}).

\begin{table}[t]
\centering
\caption{Partial Spearman correlation of twelve embedding metrics with
cross-technology transfer F1 ($n = 104$ contender conditions, 25
tissues). Controlling for embedding dimensionality $d$; $p$-values from
tissue-stratified permutation (10,000 iterations). Metrics ranked by
$|\rho|$. Bold: $p < 0.001$.}
\label{tab:metrics}
\small
\begin{tabular}{@{}lccl@{}}
\toprule
Metric & Partial $\rho$ & Perm.\ $p$ & Family \\
\midrule
Procrustes similarity   & $+0.652$ & $\mathbf{< 10^{-4}}$ & centroid \\
Silhouette (target)     & $+0.647$ & $\mathbf{< 10^{-4}}$ & scIB bio \\
Silhouette (label)      & $+0.599$ & $\mathbf{< 10^{-4}}$ & scIB bio \\
Graph connectivity      & $+0.579$ & $\mathbf{< 10^{-4}}$ & scIB bio \\
NMI (Leiden)            & $+0.553$ & $\mathbf{< 10^{-4}}$ & scIB bio \\
Cell-type CKA           & $+0.530$ & $\mathbf{< 10^{-4}}$ & centroid \\
kNN purity              & $+0.515$ & $\mathbf{< 10^{-4}}$ & scIB bio \\
cLISI                   & $+0.493$ & $\mathbf{< 10^{-4}}$ & scIB bio \\
Participation ratio     & $+0.450$ & $\mathbf{< 10^{-4}}$ & geometric \\
Isolated label ASW      & $+0.425$ & $\mathbf{< 10^{-4}}$ & scIB bio \\
ARI (Leiden)            & $+0.368$ & $\mathbf{< 10^{-4}}$ & scIB bio \\
Direction stability     & $+0.022$ & $0.83$               & geometric \\
\bottomrule
\end{tabular}
\end{table}

Controlling for both dimensionality and the number of shared cell types
attenuates correlations by $\leq 3\%$ (CKA: $\rho = 0.52$; Procrustes:
$\rho = 0.64$), confirming that the predictions are driven by embedding
quality rather than panel structure.


\subsection{The random centroid null floor}
\label{sec:nullfloor}

Predictive validity does not imply discriminative validity. We
constructed a null baseline by generating pairs of random Gaussian
centroid matrices ($k \times d$, entries $\sim \mathcal{N}(0,1)$, $d
= 512$) and computing CKA and Procrustes similarity at each $k$
matching the number of shared cell types in the panel (range $k =
8$--$44$).

Random centroid matrices achieve CKA scores of 0.92--0.99 and
Procrustes scores of 0.96--0.99 (Table~\ref{tab:nullfloor}). Both
decrease with $k$ as expected from random-matrix theory: with more cell
types, the centroid matrices sample more directions and orthogonality is
diluted. All 104 contender conditions score below their $k$-matched
random null on both metrics (contender CKA: mean 0.74, range
0.21--0.95; contender Procrustes: mean 0.78, range 0.35--0.96).

\begin{table}[t]
\centering
\caption{Random Gaussian centroid null floor. Mean CKA and Procrustes
similarity between pairs of $k \times 512$ random Gaussian matrices
(1,000 draws per $k$), with 95\% CI. $k$ = number of shared cell
types. All 104 real embedding conditions score below these baselines
on both metrics.}
\label{tab:nullfloor}
\small
\begin{tabular}{@{}rcccc@{}}
\toprule
$k$ & CKA mean & CKA 95\% CI & Procrustes mean & Procrustes 95\% CI \\
\midrule
8  & 0.985 & [0.977, 0.992] & 0.993 & [0.988, 0.996] \\
10 & 0.981 & [0.973, 0.988] & 0.991 & [0.986, 0.994] \\
15 & 0.972 & [0.964, 0.979] & 0.986 & [0.982, 0.990] \\
20 & 0.963 & [0.955, 0.970] & 0.981 & [0.977, 0.984] \\
27 & 0.950 & [0.942, 0.957] & 0.974 & [0.970, 0.977] \\
44 & 0.921 & [0.914, 0.928] & 0.958 & [0.954, 0.962] \\
\bottomrule
\end{tabular}
\end{table}

The mechanism is near-orthogonality of high-dimensional random vectors.
For $k$-dimensional subspaces of $\R^d$ with $k \ll d$, random linear
subspaces are nearly orthogonal with high probability. CKA computed on
$k \times d$ matrices is dominated by the trace inner product of the
Gram matrices; when columns are independent Gaussian, the expected CKA
approaches 1 as $d/k \rightarrow \infty$. Embeddings that preserve
biological structure reduce this independence, producing \emph{lower}
CKA than random.

Null-normalizing each condition's CKA and Procrustes scores by
subtracting the $k$-matched random mean and dividing by the random
standard deviation preserves predictive validity (normalized CKA:
$\rho = +0.53$; normalized Procrustes: $\rho = +0.66$;
$< 1\%$ change). The null floor is a scale artifact: it shifts all
absolute scores but does not distort relative rankings. A practitioner
using CKA to select among trained models would make the same choices
with or without null normalization, but would reach a wrong conclusion
if interpreting an absolute CKA of 0.74 as evidence of high
centroid alignment.


\subsection{Bio-conservation metrics under controlled degradation}
\label{sec:noise}

A metric that tracks embedding quality should decrease monotonically
as the embedding is degraded. We tested this by adding Gaussian noise at
seven levels ($\sigma \in \{0.01, 0.05, 0.1, 0.25, 0.5, 1.0, 2.0\}
\times$ embedding standard deviation) to six embedding methods across
four tissues (24 conditions).

Every bio-conservation metric responds non-monotonically
(Table~\ref{tab:noise}). ARI is non-monotonic in 22 of 24 conditions;
NMI in 17; graph connectivity in 22; cLISI in 11. At intermediate noise
levels, scores \emph{increase} before collapsing at high noise. The
mechanism differs by metric: noise drives cells into fewer Leiden
clusters, mechanically raising NMI; noise smooths local neighborhoods,
transiently increasing graph connectivity. Both batch-correction metrics
with computable values (silhouette-batch, PCR comparison) pass the
monotonicity test (23/24 each), but for a mechanical reason: noise
destroys batch structure, so batch-mixing scores improve regardless of
embedding quality.

\begin{table}[t]
\centering
\caption{Noise dose-response. Non-monotonicity count across 24
conditions (4 tissues $\times$ 6 embeddings) at 7 noise levels. A
metric passes if $\leq 1$ condition is non-monotonic. All
bio-conservation metrics fail.}
\label{tab:noise}
\small
\begin{tabular}{@{}lccc@{}}
\toprule
Metric & Non-monotonic & Total & Verdict \\
\midrule
ARI (Leiden)         & 22 & 24 & FAIL \\
Graph connectivity   & 22 & 24 & FAIL \\
NMI (Leiden)         & 17 & 24 & FAIL \\
cLISI                & 11 & 24 & FAIL \\
Isolated label ASW   &  9 & 24 & FAIL \\
Silhouette (label)   &  3 & 24 & FAIL \\
iLISI                &  2 & 24 & FAIL \\
\midrule
Silhouette (batch)   &  1 & 24 & PASS \\
PCR comparison       &  1 & 24 & PASS \\
\bottomrule
\end{tabular}
\end{table}

Non-monotonicity means that a metric can report improvement when an
embedding is degraded. Combined with the null-floor result, the pattern
is consistent: metrics built from pairwise distances, neighbor graphs,
and cluster geometry respond to geometric properties of the embedding
space that do not track biological content monotonically.


\subsection{External validation}
\label{sec:external}

Cell-type CKA was tested on three datasets with no overlap with the
Census evaluation: pancreas \citep{luecken2022scib} (6 technologies, 56
assay pairs), Tabula Sapiens \citep{tabulasapiens2022} (10x and
Smart-seq2 across 8 organs, 16 pairs), and PBMC \citep{ding2020systematic}
(7 technologies, 28 pairs). Embedding models are PCA-only baselines (no
foundation models), isolating metric behavior from model choice.

CKA generalizes to two of three datasets
(Table~\ref{tab:external}). In pancreas, CKA achieves partial
$\rho = +0.59$ ($p = 0.0001$); in Tabula Sapiens, $\rho = +0.60$
($p = 0.035$). Procrustes similarity is consistently positive across all
three ($\rho = +0.56$, $+0.73$, $+0.11$), matching the Census direction
but reaching significance in none---a power limitation from fewer
technology pairs per dataset. PBMC shows a ceiling effect: all seven
technologies produce CKA $> 0.92$ on canonical cell types (T cells, B
cells, monocytes), leaving insufficient variance for discriminative
ranking ($\rho = +0.09$, $p = 0.87$; Supplementary~\ref{app:pbmc}). CKA
predicts transfer when cell-type difficulty varies across technologies
and fails when the task is uniformly easy.

\begin{table}[t]
\centering
\caption{External validation on three independent datasets (no Census
data, no foundation models). Partial Spearman $\rho$ controlling for
embedding dimensionality; tissue-stratified permutation $p$-values.
Bold: $p < 0.05$.}
\label{tab:external}
\small
\begin{tabular}{@{}lcccccc@{}}
\toprule
& \multicolumn{2}{c}{Pancreas} & \multicolumn{2}{c}{Tabula Sapiens} & \multicolumn{2}{c}{PBMC} \\
\cmidrule(lr){2-3}\cmidrule(lr){4-5}\cmidrule(lr){6-7}
Metric & $\rho$ & $p$ & $\rho$ & $p$ & $\rho$ & $p$ \\
\midrule
Cell-type CKA & $+0.592$ & $\mathbf{0.0001}$ & $+0.604$ & $\mathbf{0.035}$ & $+0.089$ & 0.874 \\
Procrustes    & $+0.563$ & 0.115              & $+0.726$ & 0.218             & $+0.113$ & 0.659 \\
\midrule
Silhouette (tgt) & $+0.172$ & $\mathbf{<0.001}$ & $+0.664$ & $\mathbf{0.032}$ & $+0.627$ & $\mathbf{<0.001}$ \\
MMD              & $-0.451$ & 1.000              & $-0.275$ & 0.744             & $-0.135$ & 1.000 \\
Domain AUC       & $-0.495$ & 0.095              & $+0.015$ & 0.839             & $-0.148$ & 0.801 \\
\bottomrule
\end{tabular}
\end{table}


\section{Discussion}
\label{sec:discussion}

Predictive validity and discriminative validity are orthogonal
properties of embedding evaluation metrics. A metric can rank trained
models correctly by downstream transfer while scoring every trained
model below unstructured random noise. We demonstrate this dissociation
empirically: CKA and Procrustes on cell-type centroids predict
cross-technology transfer ($\rho = +0.53$ and $+0.66$) while placing
all 104 conditions below random Gaussian baselines.

The null floor arises from the geometry of high-dimensional random
matrices. Random $k \times d$ matrices with $k \ll d$ produce
near-orthogonal column subspaces, yielding high CKA and Procrustes
scores between any two independent draws. Real embeddings introduce
correlations among cell-type centroids---the biological signal that
drives transfer prediction---and these correlations \emph{reduce}
similarity scores relative to the independent-column baseline. The
paradox resolves: high CKA between random matrices reflects geometric
regularity (independence), while lower CKA between real embeddings
reflects biological structure (dependence). The metric detects the
structure that matters for transfer precisely because it deviates from
the random baseline.

The same geometric confound appears in neural population geometry, where
19 metrics are confounded by neuron count \citep{tower2026bracket}, and
is consistent with Johnson--Lindenstrauss distance preservation
\citep{johnsonlindenstrauss1984}: metrics built from pairwise distances
and neighbor graphs can remain stable even when task-relevant features
are absent. The recurrence across biological domains suggests that
geometric evaluation in high dimensions requires careful null-model
calibration whenever $d \gg k$.

The noise dose-response reveals a complementary failure mode.
Bio-conservation metrics respond non-monotonically to controlled
degradation---scores increase at intermediate noise before
collapsing---showing that these metrics do not track embedding quality
on a monotone scale. A practitioner comparing two embeddings at
different noise levels could observe an improving metric score on a
degrading embedding.

Procrustes similarity achieves the strongest transfer prediction in the
Census panel ($\rho = +0.66$) and is consistently positive across all
three external datasets ($\rho = +0.56$, $+0.73$, $+0.11$). CKA
achieves significance in two of three external datasets. Target-domain
silhouette is significant in all three external datasets but is likely
confounded: target assays with well-separated cell types are easier to
classify regardless of source embedding quality.

\paragraph{Practical recommendations.} First, use relative rankings
across models, controlling for embedding dimensionality; absolute metric
scores are not comparable to random baselines and do not indicate the
presence or absence of biological signal. Second, always include
matched-dimensionality null baselines (random projection, untrained
encoder) in evaluation panels. Third, test metric behavior under
controlled degradation before adopting a new evaluation metric. Fourth,
prefer Procrustes similarity over CKA for cross-technology transfer
prediction when cell-type labels are available.

\paragraph{Limitations.} All evaluations use cross-technology cell-type
transfer F1 as ground truth; alternative criteria such as perturbation
prediction may reveal different metric properties. The Census panel
spans 25 tissues and eight embedding methods across three architectural
families (transformer, VAE, PCA); expansion to additional models would
increase statistical power. The noise dose-response experiment uses four
tissues and six methods (24 conditions); non-monotonicity rates may
differ on larger panels. The PBMC ceiling effect limits external
validation to two of three datasets. We release all scripts,
pre-registration records, and result files to support replication.


\section{Methods}

\subsection{Data}

All embeddings are accessed via the CellxGene Census API (v2023-12-15)
\citep{abdulla2025cellxgene}. Cross-technology pairs compare 10x
Chromium 3$'$ v3 (EFO:0009922, source) against Smart-seq2
(EFO:0008931, target) within the same tissue, with zero donor overlap
enforced by metadata filtering. Tissues were auto-discovered by
requiring $\geq 200$ cells per assay and $\geq 8$ shared cell types,
yielding 25 tissues. All evaluations subsample $n = 2{,}000$ cells per
side.

Eight embedding models: Geneformer ($d = 512$)
\citep{theodoris2023geneformer}, Geneformer-v2 104M and 316M,
scGPT ($d = 512$) \citep{cui2024scgpt}, scVI ($d = 50$)
\citep{lopez2018scvi}, bag-of-genes PCA-512 ($d = 512$;
$\log(1 + x)$ without total-count normalization, PCA with fixed seed;
following \citealp{cole2026morpheus}), random Gaussian projection
($R \in \mathbb{R}^{512 \times G}$, $R_{ij} \sim \mathcal{N}(0,
1/\sqrt{512})$), and an untrained two-layer MLP encoder (random weights,
$d = 512$). Six contender models (excluding random projection and
untrained encoder) produce 104 contender conditions.

\paragraph{External validation datasets.} Three independent datasets
were used for external validation (Section~\ref{sec:external}), each
with PCA-only embeddings. Pancreas: the scIB benchmark dataset
\citep{luecken2022scib}, six technologies, 56 assay pairs, 224
contender conditions. Tabula Sapiens \citep{tabulasapiens2022}: 10x and
Smart-seq2 across eight organs, 16 pairs, 64 contenders. PBMC
\citep{ding2020systematic}: seven technologies, 28 pairs, 112
contenders.

\subsection{Cell-type centroid metrics}

Cell-type CKA \citep{kornblith2019similarity} computes centered kernel
alignment on $k \times d$ centroid matrices (mean embedding per shared
cell type) between source and target assays, using linear kernels.
Procrustes similarity applies orthogonal Procrustes analysis to
column-centered centroid matrices after PCA reduction to $\min(50,
k-1, d)$ dimensions, and returns $1 - d_{\mathrm{Proc}}^2$ after
optimal rotation. Both metrics require cell-type labels and are
evaluated on the $k$ cell types shared between source and target.

\subsection{Random centroid null}

For each value of $k$ observed in the panel ($k = 8, 9, \ldots, 44$),
we generate 1,000 pairs of random Gaussian matrices ($k \times 512$,
entries $\sim \mathcal{N}(0,1)$) and compute CKA and Procrustes
similarity. Each real condition is compared to its $k$-matched null
distribution. Null-normalized scores are computed as
$(x - \mu_k) / \sigma_k$ where $\mu_k$ and $\sigma_k$ are the null
mean and standard deviation at the condition's $k$.

\subsection{Statistical methods}

All correlations are Spearman rank with $10{,}000$-resample bootstrap
95\% CIs. Partial Spearman correlations control for embedding
dimensionality via rank residuals from OLS on ranked variables.
Tissue-stratified permutation tests (10,000 iterations) permute F1
values within tissues independently and report the proportion of
permutations with $|\rho_{\mathrm{perm}}| \geq |\rho_{\mathrm{obs}}|$.

\subsection{Pre-registration}

Scorer source code, hyperparameters, dataset specifications, and
hypothesis thresholds were frozen via SHA-256 hashing before data
access. SHA-256 hashes and the full deviation log are in the online
repository.


\section*{Data Availability}

Scorer source code, pre-registration records (SHA-256 hashes),
experiment scripts, and raw result JSON files are available at
\url{https://github.com/elliottower/preflight-bio}. All embedding data
are accessed via the CellxGene Census API.

\section*{Code Availability}

All analysis scripts are available at the repository listed under Data
Availability.

\section*{Declarations}

\paragraph{Competing interests.} The author declares no competing
interests.

\paragraph{Funding.} None. This work was conducted independently.

\paragraph{Authors' contributions.} E.T.\ conceived the study,
designed and pre-registered the protocol, performed all analyses, and
wrote the manuscript.

\paragraph{Acknowledgements.} AI tools (Claude Code) were used for code
drafting and manuscript drafting. All experimental design and conclusions
are the author's.


\appendix

\section{PBMC ceiling effect}
\label{app:pbmc}

The PBMC dataset \citep{ding2020systematic} provides seven scRNA-seq
technologies applied to the same donor PBMCs. Cell-type CKA ranges from
0.92 to 0.97 across all 28 source-target technology pairs. This
variance compression explains the null correlation with transfer F1
($\rho = +0.09$, $p = 0.87$): the metric cannot rank technologies on
a property where they score nearly identically. T cells, B cells, NK
cells, and monocytes are well-separated in any reasonable embedding; no
technology fails to capture this structure. The pancreas dataset contains
rare endocrine subtypes that are harder to resolve, providing sufficient
CKA variance for the correlation to emerge.


\section{Geometric modules}
\label{app:modules}

Five geometric modules for evaluating embedding transportability were
pre-registered and evaluated alongside the main metrics. Four were
eliminated by construct-validity checks: M1 (subspace overlap) and M6
(graph curvature) are confounded by embedding dimensionality via
Johnson--Lindenstrauss distance preservation; M2 (domain separability)
anti-predicts transfer; M4 (participation ratio) has
dimension-dependent normalization artifacts. M3 (direction stability)
was the sole candidate to survive construct-validity screening, but
achieves $\rho = +0.02$ ($p = 0.83$) on the full panel---no
predictive validity. These results demonstrate that the validity
protocol is non-trivial: it eliminates metrics, including those designed
by the protocol's authors.


\section{Noise dose-response: full data}
\label{app:noise}

Table~\ref{tab:app_noise_fail} shows representative non-monotonic
trajectories. NMI in brain/Geneformer shows scores \emph{increasing} at
high noise ($\sigma \geq 1.0$): noise drives all cells into fewer
Leiden clusters, mechanically raising NMI.

\begin{table}[h!]
\centering
\caption{Representative non-monotonic trajectories. Values are metric
scores at $\sigma \in \{0.01, 0.05, 0.1, 0.25, 0.5, 1.0, 2.0\} \times
\text{embedding\_std}$.}
\label{tab:app_noise_fail}
\footnotesize
\begin{tabular}{@{}llp{0.6\textwidth}@{}}
\toprule
Metric & Condition & $\sigma$: 0.01, 0.05, 0.1, 0.25, 0.5, 1.0, 2.0 \\
\midrule
Graph conn. & brain / scgpt & 0.772, 0.771, 0.773, \textbf{0.777}, \textbf{0.793}, 0.748, 0.589 \\
NMI & brain / geneformer & 0.550, 0.550, 0.545, \textbf{0.547}, 0.545, \textbf{0.575}, \textbf{0.577} \\
ARI & kidney / scvi & 0.397, \textbf{0.410}, 0.394, 0.351, 0.224, 0.073, 0.012 \\
cLISI & lung / scvi & 0.999, 0.999, \textbf{0.999}, 0.998, 0.975, 0.871, 0.642 \\
\bottomrule
\end{tabular}
\end{table}


\bibliographystyle{plainnat}
\bibliography{paper_d_v11}

\end{document}
